\chapter{System Design}
\label{ch:design}

\section{Architecture}

Figure~\ref{fig:arch} shows the entities and the links between them. Every
link except registration is a mutually authenticated ML-KEM/ML-DSA session
carrying AES-256-GCM records.

\begin{figure}[h]
\centering
\begin{tikzpicture}[
  box/.style={draw, rounded corners=2pt, minimum height=9mm, minimum width=20mm, align=center, font=\small, fill=white},
  lab/.style={font=\scriptsize, align=center, fill=white, inner sep=1.5pt},
  layer/.style={draw, dashed, rounded corners=4pt, inner sep=7pt, gray},
  >={Stealth[length=2.2mm]}]
\node[box] (cav1) at (0,1.2) {CAV};
\node[box] (cav2) at (0,-1.2) {CAV};
\node[box] (rsu1) at (3.6,1.2) {RSU};
\node[box] (rsu2) at (3.6,-1.2) {RSU};
\node[box] (mec) at (6.6,0) {MEC};
\node[box] (cons) at (10.2,1.4) {Consensus nodes\\(AGS-PBFT)};
\node[box] (cloud) at (10.2,-1.4) {Cloud archive};
\node[box] (ta) at (6.6,3.1) {Trusted Authority};
\node[layer, fit=(cav1)(cav2), label={[font=\scriptsize]below:user layer}] {};
\node[layer, fit=(rsu1)(rsu2)(mec), label={[font=\scriptsize]below:edge layer}] {};
\draw[<->] (cav1) -- node[lab, above=1pt]{802.11p} (rsu1);
\draw[<->] (cav2) -- node[lab, below=1pt]{802.11p} (rsu2);
\draw[<->] (rsu1) -- node[lab]{V2V relay} (rsu2);
\draw[<->] (rsu1) -- (mec);
\draw[<->] (rsu2) -- (mec);
\draw[->] (mec) -- node[lab, above, sloped]{critical} (cons);
\draw[->] (mec) -- node[lab, below, sloped]{routine} (cloud);
\draw[dashed, gray] (ta.west) -- (3.6,3.1);
\draw[->, dashed, gray] (3.6,3.1) -- node[lab, left=2pt, text=gray]{certificates,\\CRL} (rsu1.north);
\end{tikzpicture}
\caption{Entities and links. The MEC sends alerts back through the RSUs before it submits a critical event to consensus.}
\label{fig:arch}
\end{figure}

\begin{description}
\item[Trusted Authority (TA).] Generates the root key pair, certifies the
public keys of all entities, keeps the registry database and signs the
certificate revocation list (CRL).
\item[CAV.] A vehicle. It holds its own ML-DSA key pair and certificate,
opens a session with the nearest RSU, sends beacons at 10\,Hz and reports
events.
\item[RSU.] Terminates vehicle sessions, relays V2V messages to neighbouring
RSUs, and forwards events to the MEC over the wired backhaul.
\item[MEC.] The edge server. It routes events by criticality: a critical
event is first turned into an alert sent to every RSU (and so to every
vehicle), \emph{then} submitted to consensus; routine data goes to the cloud.
\item[Consensus nodes.] Eight validators running AGS-PBFT. Each committed
block carries the event, the agreed validation result, a commit certificate,
the score table and the consensus membership.
\item[Cloud.] An archive for routine data.
\end{description}

\section{Threat model}

The adversary controls the radio: it can record, drop, delay, replay, modify
and inject any frame, and it can run its own transmitters. It may obtain copies
of any certificate (they are public), run a rogue certificate authority, and
hold credentials that the TA has since revoked. Up to $f$ consensus nodes may
be Byzantine (arbitrary behaviour, including colluding on a wrong result), and
a leader may crash. In the long term the adversary has a quantum computer and
replays recorded traffic to it. The TA's private key and the private keys of
honest entities are not compromised.

\section{Registration and PKI}

\begin{enumerate}
\item The entity generates its own ML-DSA key pair. The TA never sees a
private key (this removes the key escrow of defect D4).
\item The entity sends $\{id, role, pk, alg\}$ and a proof of possession: a
signature over the request with the new private key. Without it an attacker
could have someone else's public key certified under its own name.
\item The TA checks the identity against its roster (the identities vetted out
of band, such as manufacturer records and the RSU deployment list), the
algorithm, the key length and the proof, and returns
\[ Cert = \mathsf{Sign}(sk_{TA},\ \{id, role, pk, alg, serial, nb, na\}). \]
\item The entity verifies the certificate with $pk_{TA}$ before using it.
\end{enumerate}

The certificate binds the \emph{role}, so a vehicle cannot present itself as
an RSU, and a validity period and serial number, so it can expire and be
revoked. An ML-DSA-65 certificate is \V{deck-cert} bytes. Revocation is a
TA-signed CRL holding a version number, an issue time and the revoked serial
numbers; nodes accept only a
newer, validly signed CRL and close any open session with a revoked peer.

\section{Handshake}
\label{sec:handshake}

Figure~\ref{fig:hs} shows the four-flight handshake between a client
(initiator, usually a vehicle) and a server (responder, usually an RSU).

\begin{figure}[h]
\centering
\begin{adjustbox}{max width=\linewidth}
\begin{tikzpicture}[font=\small, >={Stealth[length=2.2mm]}]
\def\L{0} \def\R{10.5}
\node[font=\bfseries] at (\L,0.4) {Client (CAV)};
\node[font=\bfseries] at (\R,0.4) {Server (RSU)};
\draw[gray] (\L,0) -- (\L,-8.2); \draw[gray] (\R,0) -- (\R,-8.2);
\draw[->] (\L,-0.6) -- node[above]{1. ClientHello $\{v, suite, Cert_C, n_C, ts\}$} (\R,-1.0);
\node[align=left, font=\scriptsize, anchor=west] at (\R+0.15,-1.6)
  {verify $Cert_C$ (TA sig, role, CRL)\\freshness, replay cache\\$(pk_{kem}, sk_{kem}) \leftarrow \mathsf{KeyGen}()$};
\draw[<-] (\L,-2.6) -- node[above]{2. ServerHello $\{Cert_S, pk_{kem}, n_S, \sigma_S\}$} (\R,-2.2);
\node[align=left, font=\scriptsize, anchor=east] at (\L-0.15,-3.4)
  {verify $Cert_S$, $\sigma_S$\\$(ss, c) \leftarrow \mathsf{Encaps}(pk_{kem})$\\$\sigma_C$; keys $\leftarrow$ HKDF$(ss, T)$};
\draw[->] (\L,-4.4) -- node[above]{3. ClientKey $\{c, \sigma_C, \mathit{Fin}_C\}$} (\R,-4.8);
\node[align=left, font=\scriptsize, anchor=west] at (\R+0.15,-5.5)
  {verify $\sigma_C$; $ss \leftarrow \mathsf{Decaps}(sk_{kem}, c)$\\erase $sk_{kem}$; keys; check $\mathit{Fin}_C$};
\draw[<-] (\L,-6.6) -- node[above]{4. ServerFinished $\{\mathit{Fin}_S\}$} (\R,-6.2);
\draw[<->, thick] (\L,-7.6) -- node[above]{AES-256-GCM records, counter from 1} (\R,-7.6);
\end{tikzpicture}
\end{adjustbox}
\caption{The mutually authenticated handshake. $T$ is the transcript of all four flights.}
\label{fig:hs}
\end{figure}

\begin{align*}
\sigma_S &= \mathsf{Sign}(sk_S,\ \text{``sh''} \parallel H(CH, SH_{body})) \\
\sigma_C &= \mathsf{Sign}(sk_C,\ \text{``ck''} \parallel H(CH \parallel SH \parallel CK_{body})) \\
T_{hash} &= H(CH \parallel SH \parallel CK_{body} \parallel \sigma_C) \\
(k_{c2s}, v_{c2s}, k_{s2c}, v_{s2c}) &= \mathrm{HKDF}(ss, T_{hash}) \\
\mathit{Fin}_C &= \mathrm{AES\text{-}GCM}(k_{c2s}, v_{c2s} \oplus 0, H(\text{``client\_finished''}, T_{hash}))
\end{align*}

Design properties:
\begin{itemize}
\item \textbf{Forward secrecy.} The server's KEM key pair is ephemeral and
erased after decapsulation, so a later compromise of long-term signing keys
does not reveal past session keys.
\item \textbf{Mutual authentication.} $\sigma_S$ binds the server's identity
to the client's fresh nonce and to $pk_{kem}$; $\sigma_C$ binds the client's
identity to the whole transcript including the ciphertext $c$.
\item \textbf{Key confirmation.} Because ML-KEM rejects implicitly, a
tampered $c$ silently produces a different key. The Finished messages detect
it: each side proves it derived the same keys.
\item \textbf{Downgrade protection.} The suite is in the ClientHello, must
match exactly, and is covered by both signatures.
\item \textbf{Uniform abort.} Every failed check raises the same error with no
detail, so a probing attacker cannot tell which check failed.
\item \textbf{Sans-IO.} The handshake state machines consume and return bytes
and never touch a socket, so the same code runs over TCP, in memory for tests
and benchmarks, and supplies exact message sizes to ns-3.
\end{itemize}

\section{Certificate caching}
\label{sec:cache}

With ML-DSA-65 each certificate is \V{deck-cert} bytes, and the full handshake
carries two of them: about half of its \V{deck-hs-bytes} bytes. Peers usually
hold these certificates already. RSU certificates are few and long-lived, so
they can be distributed to vehicles in advance with the CRL. A vehicle's
certificate can be kept, after its first handshake, in an edge directory at the
MEC that every RSU reaches over the wired backhaul.

When a peer is known to hold a certificate, the handshake sends a 32-byte
reference $\mathit{cref} = H(\text{``qrcav-cert''}, Cert)$ instead, in the
spirit of TLS cached information~\cite{rfc7924}:

\begin{enumerate}
\item ClientHello carries $\mathit{cref}_C$ instead of $Cert_C$ once the
vehicle's certificate is in the directory, and $known = H(Cert_S)$ when the
vehicle holds the RSU's certificate.
\item If $known$ matches, ServerHello carries $\mathit{cref}_S$ instead of
$Cert_S$.
\item If the RSU cannot resolve $\mathit{cref}_C$ (locally or from the
directory) it answers \emph{HelloRetry} $\{n_C\}$, and the client resends its
ClientHello with the full certificate and a new nonce.
\end{enumerate}

The security argument:
\begin{itemize}
\item \textbf{Checks are not cached.} A cached certificate is verified in full
on every handshake (TA signature, validity, role, expected identity, current
CRL), so revocation still takes effect at the next CRL refresh.
\item \textbf{The reference is bound by both signatures.} $\mathit{cref}$ is
a collision-resistant hash of the exact certificate bytes and lies inside the
messages that $\sigma_C$ and $\sigma_S$ cover. Claiming another vehicle's
cached certificate fails $\sigma_C$; replacing the server's reference fails
$\sigma_S$.
\item \textbf{The directory is not trusted.} A fetched certificate must hash
to the requested reference and is verified against $pk_{TA}$.
\item \textbf{The cache cannot be filled by strangers.} A server caches a
client certificate only after a handshake with its owner has completed, and
the cache is a bounded LRU.
\item \textbf{HelloRetry is unauthenticated, and that is acceptable.} It must
echo the client's nonce and is honoured once per handshake. Forging it only
forces the full-certificate path, the same property as TLS~1.3's
HelloRetryRequest~\cite{tls13}.
\end{itemize}

For DECK the handshake shrinks from \V{deck-hs-bytes} to
\V{deck-cached-bytes} bytes (\V{deck-cache-saving}\,\% smaller) and from
\V{deck-hs-frames} to \V{deck-cached-frames} radio frames. The cached
ClientHello is \V{deck-ch-cached} bytes, one frame.

\section{Record layer}

\begin{align*}
record &= ctr\ (8\,\mathrm{B}) \parallel \mathrm{AES\text{-}256\text{-}GCM}(k_{dir},\ nonce,\ pt,\ aad = dir \parallel ctr) \\
nonce &= v_{dir} \oplus (0^{32} \parallel \mathrm{be64}(ctr))
\end{align*}

Counter 0 is reserved for the Finished message and application data starts at
1, so no (key, nonce) pair is ever reused. A 64-record sliding window accepts
reordered records once and rejects replays; it is checked \emph{before}
decryption, so a replay flood costs a lookup rather than an AES operation. A
session refuses to send more than $2^{32}$ records.

\section{Defects found and corrected}
\label{sec:defects}

\subsection{In the base paper}

\begin{table}[h]
\centering
\caption{Defects in the base paper and their corrections}
\label{tab:defects}
\begin{tabular}{@{}p{0.6cm}p{6.2cm}p{6.8cm}@{}}
\toprule
& Defect & Correction \\
\midrule
D1 & Consensus messages are signed and verified ``using Kyber'', a KEM with no signing operation. & ML-DSA signatures on every consensus message. \\
D2 & The encapsulation listing defines $k_i$ in terms of itself, and decapsulation uses $k_i$, which the receiver does not have. & The standard \textsf{Encaps}/\textsf{Decaps} interface. \\
D3 & The session key is $H(k_i \parallel k'_j)$, a value hashed with itself, since a correct KEM gives $k_i = k'_j$. & HKDF over the shared secret and the hash of the whole transcript. \\
D4 & The TA generates entities' private keys (key escrow) although the paper claims ``no key escrow''. & Entities generate their own keys; the TA certifies public keys only, with proof of possession. \\
D5 & Hyperledger Fabric 2.2 orders with Raft and cannot host a PBFT-family protocol. & A ledger implemented directly, driven by AGS-PBFT. \\
\bottomrule
\end{tabular}
\end{table}

\subsection{In our Review 1 design}

\begin{table}[h]
\centering
\caption{Corrections to the Review 1 handshake design}
\label{tab:fixes}
\begin{tabular}{@{}p{0.6cm}p{6.2cm}p{6.8cm}@{}}
\toprule
& Review 1 & Now, and why \\
\midrule
F1 & The server signs $T \parallel pk_{kem}$, but nothing initialises $T$. & A ClientHello with a fresh nonce and timestamp opens $T$, so a recorded ServerHello cannot be replayed. \\
F2 & Only the server signs. & The client signs the transcript. Otherwise anyone holding a copy of a vehicle's (public) certificate can open a session as that vehicle. \\
F3 & Finished uses counter 0; where data starts is unspecified. & Data starts at counter 1: no GCM nonce reuse. \\
F4 & Key generation shown as $sk_{kem} \leftarrow (s, e, pk_{kem})$. & \textsf{KeyGen()}; FIPS 203's $dk$ is $(s, ek, H(ek), z)$. \\
F5 & The certificate signs only $ID \parallel pk$. & Role, serial and validity are bound too. \\
F6 & No downgrade protection. & Suite in the ClientHello, matched exactly, signed. \\
F7 & Abort reasons unspecified. & One uniform error. \\
\bottomrule
\end{tabular}
\end{table}

\section{Edge routing}

Each event carries a criticality level set by the reporting vehicle, and the
MEC routes on it. A critical event (collision ahead, pedestrian on the road,
ice) is broadcast as an alert to every RSU, which re-encrypts it
for each vehicle's session, and only then submitted to consensus for the
ledger. A routine event (a pothole report, road condition) goes to the cloud
archive. Safety warnings therefore never wait for a consensus round.

\section{AGS-PBFT}

\subsection{Normal case}

\begin{enumerate}
\item The MEC sends a signed REQUEST to the leader, $R[v \bmod |R|]$, where
$R$ is the sorted consensus set.
\item The leader sends a signed PRE-PREPARE to the consensus set and to the
candidates.
\item Consensus nodes broadcast PREPARE. On $2f+1$ matching prepares each node
\emph{validates} the event (fields present, reporter's TA certificate valid,
event not already recorded) and sends its result in a RESPONSE to the leader.
Candidates send a SHADOW-RESPONSE that is scored but not counted towards the
quorum.
\item On $2f+1$ matching responses the leader builds a block containing the
request, the agreed result, the signed responses as a commit certificate, the
updated scores and the membership, signs it, sends COMMIT to every node and
REPLY to the MEC.
\item Every node re-verifies the block (hash link, leader signature, commit
certificate, recomputed scores and membership) before appending it.
\end{enumerate}

\subsection{Scores and regrouping}

A node gains $+1$ if its result matches the committed one and loses $5$ if it
differs. Every \emph{regroup interval} (50 blocks in the paper), nodes compute
$\mu$ and $\sigma$ over all scores; consensus nodes below $\mu - \sigma$ are
demoted and candidates above $\mu + \sigma$ promoted, and the set is resized
to $3f+1$. Because scores and membership are written into each block and
re-derived by every node, all nodes agree on the consensus set without extra
messages.

\subsection{Decisions where the paper is undefined}

\begin{table}[h]
\centering
\caption{Design decisions for AGS-PBFT}
\label{tab:dd}
\begin{tabular}{@{}p{0.6cm}p{6.2cm}p{6.8cm}@{}}
\toprule
& Gap in the paper & Decision \\
\midrule
D6 & All nodes start at 100, so at the first regroup $\sigma = 0$ and $\mu - \sigma = \mu + \sigma$: every node is both above and below threshold. & Skip regrouping while $\sigma = 0$. \\
D7 & Candidates never take part, so they can never gain score or be promoted. & Candidates validate as shadows and are scored. Absent nodes score 0, not $-5$, so a Byzantine leader cannot punish honest nodes by omitting their responses. \\
D8 & ``Group size is adjusted based on network activity'' is not specified. & Request rate $r$ over the interval, from request timestamps so every node computes the same value: $f$ falls linearly from $f_{max}$ at $r \le 2$/s to 1 at $r \ge 20$/s. \\
D9 & $2f+1$ prepares are required, but the leader sends none, so the quorum is unreachable with $f$ faulty backups. & The leader also sends a PREPARE for its own proposal. \\
\bottomrule
\end{tabular}
\end{table}

\subsection{View change}

A consensus node that has seen a request but not its commit within a timeout
broadcasts VIEW-CHANGE$\{v+1, h, head, P\}$, where $P$ holds its prepared
certificates above its ledger height $h$. On $2f+1$ view changes the new
leader computes the set $O$ of requests to re-issue (the highest-view prepared
request for each sequence number, or a no-op for a gap) and sends
NEW-VIEW$\{v+1, V, O\}$. Replicas verify the signatures, recompute $O$, and on
a match enter the new view. The timeout doubles after each failed view.
Messages for a future view are buffered, because a pre-prepare can overtake
the NEW-VIEW that precedes it.

\section{Ledger}

A block is $(B\_ID, Pr\_B\_Hash, T\_Block)$ as in the paper, plus view,
sequence number, signed request, agreed result, commit certificate, score
table, membership, leader identity and leader signature. Any party holding
$pk_{TA}$ can replay a ledger from genesis and check every hash link,
signature, certificate, score and membership change.
